% =========================================================
% Proof: Finite Linear Combination Rule
% Source: volume-iii/analysis/differentiation/notes/notes-algebra.tex
% Mode: standard
% =========================================================

\subsection*{Finite Linear Combination Rule}
\label{prf:finite-linear-combination-rule}

\begin{remark*}[Return]
\hyperref[thm:finite-linear-combination-rule]{$\leftarrow$ Back to Theorem (Finite Linear Combination Rule) in Notes}
\end{remark*}

\begin{proof}
Let $I \subseteq \mathbb{R}$ be an open interval, let $c \in I$, let
$n \in \mathbb{N}$, let $f_1, \dots, f_n : I \to \mathbb{R}$ be
differentiable at $c$, and let $\alpha_1, \dots, \alpha_n \in \mathbb{R}$.
Define $h := \sum_{i=1}^{n} \alpha_i f_i$. The claim is that $h$ is
differentiable at $c$ and $h'(c) = \sum_{i=1}^{n} \alpha_i f_i'(c)$.

The proof proceeds by induction on $n$.

\ProofStep{1}{Base case $n = 1$.}
$h = \alpha_1 f_1$. By \ByThm{Constant Multiple Rule}, $h$ is differentiable
at $c$ and $h'(c) = \alpha_1 f_1'(c)$. This is $\sum_{i=1}^{1} \alpha_i
f_i'(c)$, so the claim holds for $n = 1$.

\ProofStep{2}{Inductive hypothesis.}
Assume the claim holds for some $n = k \geq 1$: for any $k$ functions
differentiable at $c$ and any scalars $\alpha_1, \dots, \alpha_k \in
\mathbb{R}$,
\[
  \left(\sum_{i=1}^{k} \alpha_i f_i\right)'(c) = \sum_{i=1}^{k} \alpha_i f_i'(c).
\]

\ProofStep{3}{Inductive step: $n = k+1$.}
Write $h = \sum_{i=1}^{k+1} \alpha_i f_i = \left(\sum_{i=1}^{k} \alpha_i
f_i\right) + \alpha_{k+1} f_{k+1}$. Let $g := \sum_{i=1}^{k} \alpha_i f_i$.
By the inductive hypothesis, $g$ is differentiable at $c$ and $g'(c) =
\sum_{i=1}^{k} \alpha_i f_i'(c)$. Since $\alpha_{k+1} f_{k+1}$ is
differentiable at $c$ by \ByThm{Constant Multiple Rule} with derivative
$\alpha_{k+1} f_{k+1}'(c)$, and since $h = g + \alpha_{k+1} f_{k+1}$,
the \ByThm{Sum Rule} gives
\[
  h'(c) = g'(c) + \alpha_{k+1} f_{k+1}'(c)
  = \sum_{i=1}^{k} \alpha_i f_i'(c) + \alpha_{k+1} f_{k+1}'(c)
  = \sum_{i=1}^{k+1} \alpha_i f_i'(c).
\]

\ProofStep{4}{Conclusion.}
By induction, the claim holds for all $n \in \mathbb{N}$.
\end{proof}

\begin{remark*}[Proof shape]
Induction on $n$, with the base case handled by the Constant Multiple Rule
and the inductive step by peeling off the last term and applying both the
Constant Multiple Rule (to scale $f_{k+1}$) and the Sum Rule (to add it to
the inductive hypothesis). The Finite Sum Rule and Constant Multiple Rule are
subsumed — this theorem is their common generalisation.

\FlashProof{Induction on $n$. Base ($n=1$): Constant Multiple Rule.
  Step: write $\sum_{i=1}^{k+1} \alpha_i f_i = \bigl(\sum_{i=1}^{k}
  \alpha_i f_i\bigr) + \alpha_{k+1}f_{k+1}$. Apply inductive hypothesis
  to the first part, Constant Multiple Rule to the second, Sum Rule to
  combine. Conclude $h'(c) = \sum_{i=1}^{k+1} \alpha_i f_i'(c)$.}
\end{remark*}

\begin{remark*}[Commentary]
\textbf{(a) Why induction rather than a direct limit argument?}
A direct limit argument is possible — split the difference quotient of
$\sum \alpha_i f_i$ term by term and apply the limit laws — but it requires
justifying that the limit distributes over a finite sum, which is itself an
induction on $n$. The induction-on-$n$ structure makes the dependency
explicit: the $n$-function result genuinely requires the $(n-1)$-function
result, and the two-function result (the Sum Rule) is the inductive engine.

\textbf{(b) Why peel off the last term rather than splitting differently?}
The inductive step writes $\sum_{i=1}^{k+1} = \bigl(\sum_{i=1}^{k}\bigr) +
(\text{one more term})$. This is the standard decomposition because the
inductive hypothesis applies to exactly $k$ terms. Any other split — say,
first term plus the rest — works equally well by the same argument. The
choice of which term to peel is arbitrary.

\textbf{(c) What does this theorem say structurally?}
The Finite Linear Combination Rule is the statement that differentiation is
a \emph{linear operator} on the vector space of differentiable functions on
$I$: it respects scalar multiplication (Constant Multiple Rule) and addition
(Sum Rule). This theorem packages both properties into a single statement
over finite linear combinations, which is exactly linearity. Every result
about linearity of differentiation that appears downstream — including the
relationship between differentiation and polynomial algebra — rests on this
foundation.
\end{remark*}

\begin{remark*}[Dependencies]
\begin{itemize}
  \item \hyperref[thm:constant-multiple-rule]{Constant Multiple Rule}:
    Handles the base case ($n=1$) and licenses scaling of $f_{k+1}$ in
    the inductive step.
  \item \hyperref[thm:sum-rule]{Sum Rule}: Combines the inductive
    hypothesis with the scaled $f_{k+1}$ term in the inductive step.
\end{itemize}

\FlashUses{thm:constant-multiple-rule, thm:sum-rule}
\end{remark*}